\begin{frame}{왜 길이가 제각각인가: 빈도 $\propto$ 노름}
  \begin{itemize}
    \item 위 임베딩에서 단어 벡터의 노름: ``사자는''=2.88,
      ``맛있다''=2.40, ``고양이는''=2.39, ``사과는''=2.00,
      ``귀엽다''=1.93 --- \textbf{같은 ``단어''라도 길이가 제각각}.
    \item \textbf{기대된 결과}: 빈도가 높은 단어가 훈련에서
      \textbf{더 많이 업데이트}되어 벡터가 더 길어진다.
    \item 실전 word2vec의 후처리에서 임베딩을 \textbf{단위벡터로
      정규화}하는 것이 이 효과를 제거하는 표준 방법
      (Lin \& Mikolov, 2015).
  \end{itemize}
  \vskip0.8em
  \kq{노름은 ``빈도''가 아니라 ``의미''가 아니다 ---
  그래서 유사도 계산은 각도(코사인)로.}
\end{frame}
